% Part: sets-functions-relations
% Chapter: sets
% Section: russells-paradox

\documentclass[../../../include/open-logic-section]{subfiles}


\begin{document}

\olfileid{sfr}{set}{rus}
\olsection{रसेलची विरोधापत्ती}

विस्तारात्मकतेमुळे $\Setabs{x}{\phi(x)}$ ही चिन्हपद्धती $\phi(x)$ असणाऱ्या
सर्व $x$ च्या \emph{त्या विशिष्ट} संचासाठी वापरणे योग्य ठरते.
मात्र, विस्तारात्मकतेमुळे \emph{प्रत्यक्षात} फक्त पुढील विचार योग्य ठरतो.
ज्याचे सदस्य सर्व आणि केवळ $\phi$ असणाऱ्या वस्तू आहेत असा संच \emph{असेल},
\emph{तर} असा एकच संच असतो. दुसऱ्या शब्दांत: एखादा $\phi$ निश्चित केल्यावर,
$\Setabs{x}{\phi(x)}$ हा संच \emph{अस्तित्वात असल्यास} तो एकमेव असतो.

पण ही अट महत्त्वाची आहे!{} प्रत्येक गुणधर्मापासून \emph{गुणधर्माधारित संचरचना}
करता येतेच असे नाही. म्हणजेच, काही गुणधर्म संचांची व्याख्या करत \emph{नाहीत}.
सर्वच गुणधर्मांनी संचांची व्याख्या होत असती, तर आपल्याला थेट व्याघातांना
सामोरे जावे लागले असते. याचे सर्वांत प्रसिद्ध उदाहरण म्हणजे रसेलची विरोधापत्ती.

संच हे इतर संचांचे !!{element}s असू शकतात; उदाहरणार्थ, $A$ या संचाचा घातसंच
हा संचांनी बनलेला असतो. म्हणून एखादा संच दुसऱ्या संचाचा !!a{element} आहे का,
असा प्रश्न विचारणे किंवा त्याचा शोध घेणे अर्थपूर्ण आहे. संच स्वतःचाच सदस्य
असू शकतो का? संचाच्या कल्पनेत असे काही दिसत नाही की ज्यामुळे ही शक्यता
नाकारली जाईल. उदाहरणार्थ, \emph{सर्व} संच मिळून वस्तूंचा समूह बनत असेल,
तर त्या सर्वांना एका संचात एकत्र करता येईल असे वाटू शकते; तो म्हणजे सर्व
संचांचा संच. तो स्वतः एक संच असल्यामुळे सर्व संचांच्या संचाचा !!a{element} असेल.

स्वतःला !!a{element} म्हणून न सामावणाऱ्या, म्हणजे \emph{स्वतःचे सदस्य नसणाऱ्या}
संचांचा गुणधर्म विचारात घेतल्यावर रसेलची विरोधापत्ती उद्भवते.
स्वतःला !!a{element} म्हणून न सामावणाऱ्या सर्व संचांचा एक संच आहे असे आपण
मानले, तर काय होईल? असा
\[
R = \Setabs{x}{x \notin x}
\]
संच अस्तित्वात आहे का? तो अस्तित्वात नसतो हे आपण सिद्ध करू शकतो.

\begin{thm}[रसेलची विरोधापत्ती]\ollabel{thm:russells-paradox}
	$R = \Setabs{x}{x \notin x}$ असा कोणताही संच नाही.
\end{thm}

\begin{proof}
$R = \Setabs{x}{x \notin x}$ अस्तित्वात असेल, तर
$R \in R$ तेव्हा आणि केवळ तेव्हाच, जेव्हा $R \notin R$; हा व्याघात आहे.
\end{proof}

\begin{tagblock}{novice}
\begin{explain}
ही सिद्धता अधिक सावकाश पाहू. $R$ अस्तित्वात असेल, तर $R \in
R$ आहे की नाही हा प्रश्न अर्थपूर्ण आहे. खरोखर $R \in R$ आहे असे समजा.
आता, स्वतःचे !!{element}s नसणाऱ्या सर्व संचांचा संच अशी $R$ ची व्याख्या केली होती.
त्यामुळे $R \in R$ असेल, तर $R$ मध्ये स्वतःच $R$ चा व्याख्यक गुणधर्म नसतो.
पण हा गुणधर्म असणारेच संच $R$ मध्ये असतात. म्हणून $R$ हा $R$ चा
!!a{element} असू शकत नाही; म्हणजेच $R \notin R$.
परंतु $R$ हा $R$ चा !!a{element} आहे आणि नाही, असे दोन्ही होऊ शकत नाही;
म्हणून व्याघात मिळतो.

$R \in R$ या गृहीतकामुळे व्याघात येत असल्याने $R \notin R$ मिळते.
पण यामुळेही व्याघात येतो!{} कारण $R
\notin R$ असेल, तर $R$ मध्ये स्वतःच $R$ चा व्याख्यक गुणधर्म असतो;
म्हणून स्वतःचे सदस्य नसणाऱ्या इतर सर्व संचांप्रमाणे $R$ हा $R$ चा
!!a{element} असेल. आणि पुन्हा, $R$ चा !!a{element} नसणे आणि असणे
या दोन्ही गोष्टी एकत्र शक्य नाहीत.
\end{explain}
\end{tagblock}

\begin{digress}
रसेलची विरोधापत्ती टाळणारा, म्हणजे $R = \Setabs{x}{x \notin x}$
अस्तित्वात आहे हा \emph{विसंगत} दावा न करणारा संच सिद्धांत कसा मांडायचा?
त्यासाठी संच केव्हा अस्तित्वात असतात (आणि केव्हा नसतात) हे सांगणाऱ्या
अगदी नेमक्या अटी देणारी स्वयंसिद्धके मांडावी लागतील.

या प्रकरणात आराखडा दिलेला संच सिद्धांत असे करत नाही. तो
\emph{खरोखरच अनौपचारिक} आहे. त्यातून एवढेच सांगितले जाते की संच
विस्तारात्मकतेचे पालन करतात आणि काही संच दिलेले असतील, तर त्यांचा संयोग,
छेद इत्यादी बनवता येतात. संच सिद्धांत यापेक्षा अधिक काटेकोरपणे विकसित करणे
शक्य आहे. \oliflabeldef{cumul:::part}{तो काटेकोरपणा \olref[cumul][][]{part}
या भागासाठी राखून ठेवू. सध्या आपण अनौपचारिक रीतीने पुढे जाऊ आणि व्याघात
टाळण्याचा काळजीपूर्वक प्रयत्न करू.}{}
\end{digress}


\end{document}
